\paragraph{\trAlgebraicTwisting{}与\trGeometricTwisting。}\label{sub:algebraic-geometric-twisting}
设有一个底为\trDisk{}的\trLefschetzFibration，其\trFiber{}为 $M$，只有一个\trCriticalPoint。若\trFiber{}维数为 $2n>2$，这类\trFibration{}总容许一个\trRelative{}二次复\trVolume{}形式，它承袭自 $M$ 的那一个；而当 $2n=2$ 时，这当且仅当\trVanishingCycle{}$V$ 容许一个\trGrading{}时才会如此。记住这一点，(17d)–(17h) 的整段论证便可照旧通过，而且所得的全部符号都与\cref{thm:quasi-representative-criterion}中所预期的一致；这是一项直截了当却略显乏味的验证，我们略去不写。正如 \cref{thm:relative-brane-local-model} 那样，当 $2n>2$ 时标准\trBoundaryCondition{}不容许\trRelative\trGrading，但这不成问题，因为在那些维数下，\cref{thm:cochain-level-vanishing} 的消没论证并不依赖任何抵消。在余下的情形中，\cref{thm:relative-brane-local-model} 适用，于是若想在 $\mathrm{char}(\KK)\neq 2$ 下工作，就必须在\trVanishingCycle{}$V=L_0$ 上选取非平凡的 \trPinStructure。主要结果取自 (17h)，但如今已配上它应有的\trGrading{}与符号；陈述如下。
\begin{theorem}\label{thm:twist-tau-isomorphism-modules}
设 $V=L_0$ 是 $M$ 中一个带\trFrame{}的\trExact\trLagrange\trSphere，为其配备\trBraneStructure{}$L^{\#}_0$（在 $n=1$ 且 $\mathrm{char}(\KK)\neq 2$ 的情形，归属于它的那个 \trPinStructure{}应当取非平凡的）。再设 $L_1$、$L_2$ 是另外两个\trExact\trLagrange\trSubmanifold，带有取定的\trIsotopy{}$L_2\simeq\tau_V(L_1)$，并设它们也带有\trBraneStructure{}$L^{\#}_k$，这些\trBraneStructure{}与该\trIsotopy\trCompatible（此处给 $\tau_V$ 配的是其自然\trGrading）。把由 $L^{\#}_k$ 给出的 $\scr{F}(M)$ 中的那些\trObject{}记作 $Y_k$，把它们的\trYoneda{} \trAinfty-模记作 $\scr{Y}_k$。那么，通过对\trFibration{}$E^c\to S^c$ 与 $E^{k,r}\to S^{k,r}$ 的\trSection{}进行计数，可得 $c\in CF^0(L^{\#}_1,L^{\#}_2)$ 与 $k\colon CF^*(L^{\#}_0,L^{\#}_1)\to CF^*(L^{\#}_0,L^{\#}_2)[-1]$，使得由此得到的、如\cref{thm:quasi-representative-criterion}中那样定义的 $[t]\colon\scr{T}_{Y_0}(\scr{Y}_1)\to\scr{Y}_2$ 是 $H^0(\mathit{mod}(\scr{F}(M)))$ 中的一个\trIsomorphism。
\end{theorem}
\noqed
为了得到一个更具体的陈述，我们如今像 \eqref{eq:twist-notions-compatible} 中那样，把 \trAinfty-模换成\trTwistedComplex。我们还假定 $L_2=\tau_V(L_1)$，并省去一些更精细的细节。这就把\cref{thm:twist-tau-isomorphism-modules}变成了下面这条“$T=\tau$”的陈述：
\begin{corollary}\label{thm:twist-equals-tau}
在\trDerivedCategory{}$D\scr{F}(M)$ 中，有 $T_{L^{\#}_0}(L^{\#}_1)\cong\tau_{L_0}(L^{\#}_1)$。
\end{corollary}
\noqed
$T_{L^{\#}_0}(L^{\#}_1)$ 处于\trExactTriangle{} \eqref{eq:twisted-yoneda-triangle} 之中，而根据\trExactTriangle{}的一般性质，可以用 $\tau_{L_0}(L^{\#}_1)$ 来替换它。注意，该三角形中的某一个\trMorphism{}是 \cref{sub:twisting-yoneda-modules} 中定义的（纯粹代数的）赋值\trMorphism{}$[f]$。据\cref{thm:quasi-representative-criterion}，第二个\trMorphism{}是 $[c]$，于是我们有：
\begin{corollary}\label{thm:hf-tau-exact-triangle}
在 $D\scr{F}(M)$ 中有\trExactTriangle
\begin{equation}\label{eq:hf-tau-exact-triangle}
\begin{tikzcd}
HF^*(L^{\#}_0,L^{\#}_1)\otimes L^{\#}_0 \ar[rr,"{[f]}"] && L^{\#}_1 \ar[dl,"{[c]}"{swap}] \\
& \tau_{L_0}(L^{\#}_1). \ar[ul,"{[1]}"] &
\end{tikzcd}
\end{equation}
\end{corollary}
\noqed
\begin{remark}\label{thm:third-map-dual-evaluation}
\eqref{eq:hf-tau-exact-triangle} 中的第三个映射是
\begin{equation*}
\begin{aligned}
&\Hom_{D\scr{F}(M)}\bigl(\tau_{L_0}(L^{\#}_1),HF^*(L^{\#}_0,L^{\#}_1)\otimes L^{\#}_0[1]\bigr)\\
&\cong\bigl(HF^*(L^{\#}_0,L^{\#}_1)\otimes HF^*\bigl(\tau_{L_0}(L^{\#}_1),L^{\#}_0\bigr)\bigr)^{1}\\
&\cong\bigl(HF^*(L^{\#}_0,L^{\#}_1)\otimes HF^*(L^{\#}_1,L^{\#}_0)\bigr)^{n},
\end{aligned}
\end{equation*}
中的一个元素，其中最后一个\trIsomorphism{}用到了 \cref{thm:relative-brane-local-model} 中提到的等式 $\tau_{L_0}(L^{\#}_0)=L^{\#}_0[1-n]$。自然会期望，这就是由诸\trBasis{}的\trPoincare{}对偶对所给出的那个元素；换句话说，可以把它看作 \eqref{eq:adjoint-twist-triangle} 中所用的那种对偶\trEvaluationMap。看起来，原则上把\cref{thm:exactness-from-conditions}与 \cref{thm:additional-data-gpq} 相结合，便能证明情形确实如此；然而我们并未核验全部细节；特别是，符号问题相当棘手，并且会以非平凡的方式影响结果（与\cref{thm:massey-product-exactness}比较）。
\end{remark}
最后，我们有一个更为实质性的应用。设 $\mathbf{V}^{\#}=(V^{\#}_1,\ldots,V^{\#}_m)$ 是 $M$ 中一组带\trFrame{}的\trExact\trLagrange\trSphere，每一个都配备了\trBraneStructure（当 $n=1$ 且 $\mathrm{char}(\KK)\neq 2$ 时，我们要求所有这些 \trPinStructure{}都是非平凡的）。设 $\mathbf{W}^{\#}=\mathscr{L}_k\mathbf{V}^{\#}$ 是由它经一次初等 \trHurwitz{}移动 \eqref{eq:hurwitz-move-spheres} 得到的那一组，并带有自然诱导的\trBraneStructure（就\trGrading{}而言，这用到 $\tau_{V_k}$ 的自然\trGrading）。设 $\scr{F}^{\rightarrow}(\mathbf{V}^{\#})$、$\scr{F}^{\rightarrow}(\mathbf{W}^{\#})$ 是 $\scr{F}(M)$ 中与之相伴的\trDirected{} \trAinfty-\trSubcategory。\cref{thm:twist-equals-tau}说，$\mathbf{W}^{\#}$ 的诸\trObject{}在 $D^b(\scr{F}(M))$ 中，与对 $\mathbf{V}^{\#}$ 施行一次\trMutation{} \eqref{eq:elementary-mutation} 所得到的那一组的诸\trObject{}是\trIsomorphic{}的。由\cref{thm:directed-collections-quasi-isomorphic}，换成\trIsomorphic{}的\trObject{}不会改变\trDirected{} \trAinfty-\trCategory{}的\trQuasiisomorphism{}类。于是\cref{thm:spherical-mutation-directed}适用，从而 $Tw\scr{F}^{\rightarrow}(\mathbf{V}^{\#})$ 与 $Tw\scr{F}^{\rightarrow}(\mathbf{W}^{\#})$ 是\trQuasiequivalent{}的。一个远为初等的事实是：倘若把某一个或某几个 $V_j$ 的\trGrading{}作\trSuspension，$Tw\scr{F}^{\rightarrow}(\mathbf{V}^{\#})$ 并不改变。我们把结果总结如下：
\begin{theorem}\label{thm:twisted-directed-hurwitz-invariant}
在相差\trQuasiequivalence{}的意义下，$Tw\scr{F}^{\rightarrow}(\mathbf{V}^{\#})$ 是 $\mathbf{V}$ 的 \trHurwitz{}\trEquivalence{}类的一个\trInvariant。
\end{theorem}
\noqed
当把它应用于\trVanishingCycle{}的\trDistinguishedBasis{}时，它立即给出下述结果：
\begin{corollary}\label{thm:twisted-directed-fibration-invariant}
设 $\pi\colon E\to S$ 是一个\trLefschetzFibration，其底 $S$ 是一个\trDisk，其上带有\trRelative{}二次复\trVolume{}形式 $\eta^2_{E/S}$。取定一个基点 $\ast$，并为 $M=E_{\ast}$ 配备与之相伴的形式 $\eta^2_{M}$。设 $\boldsymbol{\gamma}$ 是始于 $\ast$ 的\trVanishingPath{}的一个\trDistinguishedBasis，而 $\mathbf{V}_{\boldsymbol{\gamma}}$ 是 $M$ 中与之相伴的\trVanishingCycle{}组。为这些\trCycle{}中的每一个配备一个\trGrading{}与一个 \trPinStructure（当 $n=1$ 且 $\mathrm{char}(\KK)\neq 2$ 时选取非平凡的 \trPinStructure），并把所得结果记作 $\mathbf{V}^{\#}_{\boldsymbol{\gamma}}$。设 $\scr{F}^{\rightarrow}(\mathbf{V}^{\#}_{\boldsymbol{\gamma}})$ 是 $\scr{F}(M)$ 中与之相伴的\trDirected{} \trAinfty-\trSubcategory。在相差\trQuasiequivalence{}的意义下，\trTwistedComplex{}的\trCategory{}$Tw\scr{F}^{\rightarrow}(\mathbf{V}^{\#}_{\boldsymbol{\gamma}})$ 既不依赖于 $\boldsymbol{\gamma}$ 的选取，也不依赖于诸\trGrading，因而是该\trLefschetzFibration{}的一个\trInvariant。从而，在相差\trTriangulated\trCategory{}之\trEquivalence{}的意义下，\trDerivedCategory{}$D\scr{F}^{\rightarrow}(\mathbf{V}^{\#}_{\boldsymbol{\gamma}})$ 亦然。
\end{corollary}
\noqed
\begin{remark}\label{thm:kontsevich-conjecture-remark}
\cref{thm:twist-equals-tau}最初是\trKontsevich{}的一个猜想。显然，人们会期望一个更强的版本成立，也就是说，由 $\tau$ 诱导的那个\trAutomorphism{}应当在 $H^0(\mathit{fun}(\scr{F}(M),\scr{F}(M)))$ 中与 $T$ \trCanonical{}\trIsomorphic。我们将不在此证明这一点；但可以期望，它将是 Wehrheim--Woodward 关于\trLagrange{}对应的正在进行的工作的一个推论。
\end{remark}
